\begin{figure}[H]
\centering

\begin{tikzpicture}[
    scale=0.65,
    transform shape,
    neuron/.style={circle, minimum size=10mm, inner sep=0pt},
    >=stealth
]

\definecolor{darkgrey}{RGB}{90,90,90}

% -------- Step bracket tuning --------
\def\StepTick{4pt}
\def\StepRise{10pt}
\def\StepGap{8pt}
\def\TickSkew{6pt}

% ======================
% Input layer (3 visible nodes)
% ======================
\node[neuron, fill=gray!40] (S1)  at (0,  2.5) {$S_{1}$};
\node[neuron, fill=gray!40] (S2)  at (0,  1.0) {$S_{2}$};
\node at (0, -0.2) {$\vdots$};
\node[neuron, fill=gray!40] (SM)  at (0, -2.2) {$S_{M}$};

% ======================
% Latent layer
% ======================
\node[neuron, fill=red!70] (Ztop) at (4,  1.9) {};
\node[above=2mm] at (Ztop.north) {$\tau_i = 1,\ldots,30$};

\node[neuron, fill=red!70] (Z1)  at (4,  0.4) {$z_{1}$};
\node                    (Zvd) at (4, -0.6) {$\vdots$};
\node[neuron, fill=red!70] (Zd)  at (4, -1.8) {$z_{\ell}$};

% ======================
% Hidden layer (10 nodes)
% ======================
\foreach \i in {1,...,10} {
    \node[neuron, fill=red!40] (H\i) at (8, {(5.5-\i)*1.15}) {$x_{1,\mathcal{G}}^{(\i)}$};
}

% ======================
% Output layer (3 nodes)
% ======================
\node[neuron, fill=red!20] (P1)  at (12,  2.5) {$P_{\tau_1}$};
\node[neuron, fill=red!20] (P2)  at (12,  1.0) {$P_{\tau_2}$};
\node at (12, -0.2) {$\vdots$};
\node[neuron, fill=red!20] (P30) at (12, -2.2) {$P_{\tau_{30}}$};

% ======================
% Reconstructed swaps
% ======================
\node[neuron, fill=gray!10] (St1) at (16,  2.5) {$\tilde{S}_{1}$};
\node[neuron, fill=gray!10] (St2) at (16,  1.0) {$\tilde{S}_{2}$};
\node at (16, -0.2) {$\vdots$};
\node[neuron, fill=gray!10] (StM) at (16, -2.2) {$\tilde{S}_{M}$};

% ======================
% Connections
% ======================
\foreach \s in {S1,S2,SM} {
    \draw[->, draw=darkgray] (\s) -- (Z1);
    \draw[->, draw=darkgray] (\s) -- (Zd);
}

\foreach \j in {1,...,10} {
    \draw[->, draw=darkgray] (Ztop) -- (H\j);
    \draw[->, draw=darkgray] (Z1)   -- (H\j);
    \draw[->, draw=darkgray] (Zd)   -- (H\j);
}

\foreach \i in {1,...,10} {
    \draw[->, draw=darkgray]         (H\i) -- (P1);
    \draw[->, dashed, draw=darkgray] (H\i) -- (P2);
    \draw[->, dotted, draw=darkgray] (H\i) -- (P30);
}

\draw[->, draw=darkgray] (P1) -- (St1);
\draw[->, draw=darkgray] (P1) -- (St2);
\draw[->, draw=darkgray] (P2) -- (St2);
\draw[->, draw=darkgray] (P1)  -- (StM);
\draw[->, draw=darkgray] (P2)  -- (StM);
\draw[->, draw=darkgray] (P30) -- (StM);

% ======================
% Brackets
% ======================

\coordinate (StepY) at ([yshift=\StepRise]H1.north);

\newcommand{\StepBracket}[3]{%
  \draw[thick] (#1) -- ++(0,-\StepTick);
  \draw[thick] (#1) -- (#2);
  \draw[thick] (#2) -- ++(0,-\StepTick);
  \node[above=4pt] at ($(#1)!0.5!(#2)$) {\textbf{#3}};
}

% Encoder bracket
\coordinate (B0) at ([xshift=\TickSkew]S1.center |- StepY);
\coordinate (B1) at ([xshift=-\TickSkew]Ztop.center |- StepY);
\coordinate (EncL) at (B0);
\coordinate (EncR) at ([xshift=-\StepGap]B1);
\StepBracket{EncL}{EncR}{Encoder}

% Decoder bracket (spans n1 to n4)
\coordinate (B2) at ([xshift=\TickSkew]Ztop.center |- StepY);
\coordinate (B3) at ([xshift=-\TickSkew]St1.center |- StepY);
\coordinate (DecL) at ([xshift=\StepGap]B2);
\coordinate (DecR) at (B3);
\StepBracket{DecL}{DecR}{Decoder}

% ======================
% Layer labels
% ======================
\large
\node at (0,  -6.2)  {Input};
\node at (4,  -6.2)  {Latent};
\node at (8,  -6.2)  {Hidden};
\node at (12, -6.2)  {Prices};
\node at (16, -6.2)  {Output};
\normalsize

\end{tikzpicture}

\caption{
Architecture of the arbitrage-free autoencoder. The encoder maps $M$ observed swap rates $S_1,\ldots,S_M$ to $\ell$ latent factors $z_1,\ldots,z_\ell$. The decoder evaluates the latent factors against each maturity $\tau_i$ in turn, passing them through a hidden layer of ten nodes to produce one bond price $P_{\tau_i}$ at a time across the tenor grid $\tau_i = 1,\ldots,30$ (solid, dashed, and dotted arrows). The reconstructed swap rates $\tilde{S}_1,\ldots,\tilde{S}_M$ are recovered from the resulting bond prices.
}

\label{fig:rolf_architecture}

\end{figure}